\chapter{Metalwork Removal}

\section*{What Is This Surgery?}
Metalwork removal is a common elective orthopedic surgical procedure involving the extraction of internal fixation devices, such as plates, screws, rods, wires, or pins, that were previously implanted to stabilize fractures, perform osteotomies, or achieve arthrodesis. This intervention is typically performed once the underlying bone has fully healed and the implant is no longer biomechanically necessary for skeletal stability. The primary indications for removal include symptomatic hardware causing pain, soft tissue irritation, or functional impairment, as well as preventing potential long-term complications or facilitating future medical imaging or surgical interventions. It is a reconstructive procedure aimed at improving patient comfort and function after the initial orthopedic treatment has achieved its goal of bone union.

\section*{Alternative Names}
\begin{itemize}
  \item Hardware removal
  \item Implant removal surgery
  \item Internal fixation device extraction
  \item Orthopedic hardware removal
  \item Plate and screw extraction
  \item Intramedullary nail removal
\end{itemize}

\section*{Relevant Surgical Anatomy}
The relevant surgical anatomy for metalwork removal is highly dependent on the original site of implantation, but general principles apply across all regions. The surgeon must have an intimate understanding of the superficial and deep anatomical layers that were traversed during the initial implantation. These layers typically include the skin, subcutaneous fat, various fascial layers, muscle bellies, tendons, ligaments, and ultimately, the bone where the implant resides.

Crucially, the proximity of neurovascular bundles to the implant site is a paramount consideration. Major nerves (e.g., radial, ulnar, median in the upper limb; sciatic, femoral, peroneal in the lower limb) and blood vessels (e.g., brachial artery, femoral artery, popliteal artery) must be meticulously identified and protected during dissection, as they are often in close relation to the bone and the implanted hardware. Scar tissue from the initial surgery can obscure these vital structures, making careful, layer-by-layer dissection along the original surgical planes essential. Specific anatomical landmarks, such as McBurney's point for appendectomy or Calot's triangle for cholecystectomy, are not directly applicable to metalwork removal, but the surgeon relies on knowledge of bony prominences, muscle origins/insertions, and fascial septa to guide their approach back to the implant.

\begin{itemize}
  \item \textbf{Skin and Subcutaneous Tissue:} The initial incision often follows the previous surgical scar, requiring careful dissection through potentially fibrotic tissue.
  \item \textbf{Fascial Layers:} Deep fascia often encases muscles and can be a plane for dissection, but also a barrier to implant access.
  \item \textbf{Muscles and Tendons:} These structures may need to be carefully retracted or split along their fibers to expose the underlying bone and hardware.
  \item \textbf{Neurovascular Structures:} Nerves and vessels are at risk of injury, especially when embedded in scar tissue or in close proximity to the implant.
  \item \textbf{Bone:} The implant is typically seated within or on the bone, which may show signs of stress shielding or remodeling around screw holes.
\end{itemize}

\section*{Surgical Principle \& Rationale}
The fundamental principle behind metalwork removal is that internal fixation devices are generally temporary implants designed to provide mechanical stability during the bone healing process. Once bone union is achieved, the implant's primary function is complete, and its continued presence can become a source of complications or discomfort. The surgical rationale for removal is to alleviate symptoms such as pain, tenderness, or bursitis caused by the implant's prominence or irritation of surrounding soft tissues, tendons, or nerves. This involves carefully re-accessing the implant site and systematically disassembling and extracting the components.

Furthermore, removal may be indicated to address implant-related infections, documented allergic reactions to the metal, mechanical failure of the device, or to prevent potential long-term issues like stress shielding, which can lead to localized osteopenia and increase the risk of refracture upon removal. In pediatric patients, removal is often necessary to prevent growth disturbances or to accommodate normal skeletal development. The technical goal of the procedure is the complete and safe extraction of all implant components, minimizing damage to surrounding tissues, preserving the integrity of the healed bone, and restoring the anatomical and functional status of the affected limb. This often involves careful dissection along the original surgical planes and the use of specialized instruments to engage and remove the various types of fixation devices.

\section*{History \& Surgical Evolution}
The concept of internal fixation for fracture management gained significant traction in the early 20th century, with pioneers like Sir Arbuthnot Lane (1905) in the UK and William O'Neill Sherman (1912) in the US developing early plate and screw systems. Initially, these implants were made of materials like steel, which were prone to corrosion and often necessitated removal due to adverse tissue reactions and high rates of infection. The mid-20th century saw a revolution with the advent of stainless steel alloys and the establishment of the AO Foundation (Arbeitsgemeinschaft f\"ur Osteosynthesefragen) in Switzerland by Maurice E. M\"uller, Martin Allg\"ower, Robert Schneider, and Hans Willenegger (1958). The AO Foundation standardized surgical techniques and implant designs, leading to widespread adoption of internal fixation and a reduction in the need for routine removal.

As implant materials improved, particularly with the introduction of titanium and its alloys in the latter half of the 20th century, the biocompatibility of hardware increased, significantly reducing the incidence of corrosion and adverse reactions. This led to an ongoing debate regarding the necessity of routine metalwork removal versus selective removal only for symptomatic cases. While some surgeons advocated for routine removal to prevent potential long-term complications, others argued against it due to the inherent risks of a second surgical procedure. Current practice generally favors selective removal based on patient symptoms and specific indications, reflecting a more nuanced understanding of implant longevity and patient-specific needs and a shift towards evidence-based decision-making.

\section*{Clinical Indications}
Metalwork removal is indicated for a variety of reasons, primarily when the implant causes symptoms or complications.

\begin{itemize}
  \item Persistent pain or discomfort directly attributable to the implant, such as localized tenderness, deep aching, or nerve irritation.
  \item Palpable prominence or soft tissue irritation, leading to bursitis, skin breakdown, or impingement on tendons or nerves, particularly in superficial locations.
  \item Documented implant-related infection, which often necessitates removal of the foreign body to eradicate the infection completely.
  \item Confirmed allergic reaction or sensitivity to the implant materials, although this is a relatively rare occurrence and requires specific diagnostic testing.
  \item Mechanical failure of the implant, such as a broken plate, screw, or rod, which may cause pain, loss of stability, or migration of fragments.
  \item Interference with joint motion or tendon gliding, limiting functional range of motion or causing snapping phenomena.
  \item Need for future medical imaging, particularly Magnetic Resonance Imaging (MRI), where metallic implants can cause significant artifact that obscures diagnostic detail.
  \item Prevention of growth disturbance in pediatric patients, especially when implants cross or are near open physes (growth plates).
  \item Patient request for removal, even if asymptomatic, due to psychological factors or personal preference, after thorough discussion of risks and benefits.
\end{itemize}

\section*{Clinical Positioning}
Metalwork removal is almost exclusively a secondary, or reconstructive, procedure. It is performed after the primary orthopedic treatment, such as fracture fixation, osteotomy, or arthrodesis, has successfully achieved its intended goal, typically bone union. The initial surgery is considered the primary intervention, aimed at stabilizing the skeletal structure and facilitating healing. Metalwork removal then addresses issues that arise from the presence of the implant itself, rather than the original skeletal pathology.

In rare circumstances, such as acute implant-related infection or early mechanical failure, metalwork removal might be performed relatively soon after the primary surgery, but its role remains secondary to the initial fixation and is considered a salvage procedure to mitigate complications. It is not a diagnostic test but a therapeutic intervention to resolve implant-related problems. Therefore, its clinical positioning is firmly in the realm of elective, secondary surgical management, with the goal of improving the patient's long-term comfort and function.

\section*{Surgical Technique \& Procedure}
The procedure for metalwork removal typically begins with the administration of anesthesia, which can be general anesthesia, regional anesthesia (such as a spinal or epidural block), or sometimes local anesthesia with sedation, depending on the implant's location, complexity of removal, and patient comorbidities. The patient is positioned on the operating table to allow optimal access to the surgical site, often mirroring the positioning used during the initial implantation surgery to facilitate re-entry along established surgical planes.

The surgeon usually re-opens the previous surgical incision, carefully dissecting through the skin and subcutaneous tissues, often encountering scar tissue and adhesions from the prior operation. Meticulous dissection is performed to identify and expose the internal fixation device, ensuring that vital neurovascular structures are protected throughout the process. Once the implant is fully visualized, specialized instruments corresponding to the specific type of hardware (e.g., specific screwdriver bits for different screw heads, plate benders, intramedullary nail extractors) are used for removal.

The general steps for metalwork removal are as follows:
\begin{itemize}
  \item \textbf{Anesthesia and Positioning:} Administer appropriate anesthesia and position the patient to optimize surgical access and safety.
  \item \textbf{Incision and Dissection:} Re-open the original surgical incision, carefully dissecting through scar tissue and soft tissues to expose the implant, protecting neurovascular structures.
  \item \textbf{Implant Exposure:} Fully expose the plate, screws, rod, or other hardware, ensuring all components are visible and accessible.
  \item \textbf{Screw Removal:} Systematically remove all screws using the correct screwdriver bit, ensuring not to strip the screw heads. For broken screws, specialized extraction tools may be required.
  \item \textbf{Plate/Rod Extraction:} Once all screws are removed, carefully lift plates from the bone. Intramedullary rods or nails are extracted using specific extraction tools, such as slap hammers or threaded extractors, often requiring a separate small incision at the entry point.
  \item \textbf{Site Inspection and Irrigation:} After all components are removed, the surgical site is thoroughly inspected for any retained fragments, bone debris, or soft tissue damage, and then irrigated with sterile saline solution.
  \item \textbf{Hemostasis and Closure:} Achieve meticulous hemostasis. The soft tissues are then meticulously closed in layers, and the skin incision is approximated, often with sutures or staples.
  \item \textbf{Dressing Application:} A sterile dressing is applied to the wound. Drains are rarely used unless there is significant dead space, concern for hematoma formation, or a history of infection.
\end{itemize}

\section*{Preoperative Workup \& Preparation}
Preparation for metalwork removal involves a comprehensive pre-operative workup to ensure patient safety and optimize surgical outcomes. This typically includes a detailed clinical history and physical examination to assess the patient's overall health, identify any comorbidities, and evaluate the specific symptoms related to the hardware. Imaging studies are crucial; plain radiographs (X-rays) are routinely obtained to confirm complete bone union, identify the exact location, type, and number of implant components, and assess for any signs of implant loosening, breakage, or stress shielding. In complex cases, for deeply buried hardware, or to evaluate surrounding soft tissues, a Computed Tomography (CT) scan or Magnetic Resonance Imaging (MRI) may be utilized.

\begin{itemize}
  \item \textbf{Medical Evaluation:} Comprehensive medical and cardiac evaluation, including an assessment of anesthetic risk, especially for patients with significant medical comorbidities.
  \item \textbf{Laboratory Tests:} Routine pre-operative laboratory tests, such as a complete blood count, electrolyte panel, renal and liver function tests, and coagulation studies, are performed.
  \item \textbf{Medication Review:} Patients are instructed to be NPO (nil per os) for a specified period before surgery. Adjustments to regular medications, particularly anticoagulants or antiplatelet agents, are made as per institutional protocols and in consultation with the prescribing physician.
  \item \textbf{Prophylactic Antibiotics:} Prophylactic antibiotics are typically administered intravenously prior to the incision to minimize the risk of surgical site infection.
  \item \textbf{Informed Consent:} Finally, informed consent is obtained, ensuring the patient fully understands the procedure, its potential benefits, risks, alternatives, and the expected recovery course.
\end{itemize}

\section*{Contraindications \& Precautions}
\textbf{Absolute Contraindications:}
\begin{itemize}
  \item Active infection at the surgical site, unless the primary goal of the surgery is to remove infected hardware and perform thorough debridement to eradicate the infection.
  \item Unstable medical comorbidities that significantly increase the risk of anesthesia or surgery, such as severe uncontrolled cardiac disease, acute respiratory failure, or recent myocardial infarction.
  \item Incomplete bone healing or non-union of the original fracture or osteotomy, as removal of the stabilizing hardware could lead to immediate refracture, collapse, or loss of correction.
\end{itemize}

\textbf{Relative Contraindications and Precautions:}
\begin{itemize}
  \item Asymptomatic hardware: If the implant is not causing any problems, the inherent risks of surgery may outweigh the potential benefits of removal.
  \item Deeply buried or difficult-to-access hardware: Implants in anatomically challenging locations (e.g., spine, pelvis, deep within joints) may pose higher risks of neurovascular injury or incomplete removal.
  \item Hardware in critical anatomical locations: Proximity to major nerves, blood vessels, or vital organs necessitates extreme caution and may significantly increase surgical risk.
  \item Significant scarring or previous complex surgeries: Extensive scar tissue can make dissection difficult, increase operative time, and elevate the risk of tissue damage.
  \item Patient refusal: The patient's informed decision to decline surgery must be respected after a thorough discussion of all options.
  \item Pediatric patients with open growth plates: Removal may be delayed unless the hardware is causing growth disturbance, to avoid potential iatrogenic injury to the physis.
  \item Expected short life expectancy: In patients with limited life expectancy, the benefits of elective removal may not justify the surgical risks and recovery period.
  \item Broken or stripped screws: These can complicate removal and may require specialized extraction techniques or potentially leave fragments behind.
\end{itemize}

\section*{Complications \& Risks}
As with any surgical procedure, metalwork removal carries inherent risks, which can be broadly categorized into general surgical complications and procedure-specific risks.

\textbf{General Surgical Complications:}
\begin{itemize}
  \item Bleeding: Hematoma formation at the surgical site, which may cause pain, swelling, and potentially require drainage.
  \item Infection: Superficial wound infection or, less commonly, deep surgical site infection, which may necessitate antibiotic treatment or further surgical debridement.
  \item Anesthesia complications: Nausea, vomiting, allergic reactions to medications, or more severe cardiac or pulmonary events, including myocardial infarction or stroke.
  \item Pain: Post-operative pain at the incision site, which is usually manageable with analgesics but can occasionally be chronic.
  \item Scarring: Formation of a new or exacerbated scar, including hypertrophic scars or keloids, particularly in susceptible individuals.
  \item Deep Vein Thrombosis (DVT) and Pulmonary Embolism (PE): Blood clot formation, especially in the lower extremities, with the risk of migration to the lungs.
\end{itemize}

\textbf{Procedure-Specific Risks:}
\begin{itemize}
  \item Refracture: A significant risk, particularly through screw holes or areas of stress shielding, especially if the bone has not fully consolidated or if activity is resumed too quickly.
  \item Nerve or vascular injury: Damage to adjacent nerves, leading to numbness, weakness, or paralysis, or injury to blood vessels, potentially causing hemorrhage, pseudoaneurysm, or ischemia.
  \item Incomplete hardware removal: A screw head may strip, a screw may break, or a fragment of the implant may be left behind, potentially requiring further surgery or causing ongoing symptoms.
  \item Persistent pain or recurrence of symptoms: The original symptoms may not resolve if the hardware was not the sole cause of discomfort, or new pain may develop from scar tissue or nerve irritation.
  \item Damage to surrounding soft tissues: Injury to muscles, tendons, or ligaments during dissection or implant extraction, potentially leading to weakness or functional deficits.
  \item Heterotopic ossification: Abnormal bone formation in soft tissues around the surgical site, which can cause pain and restrict joint motion.
  \item Difficulty locating or extracting hardware: Especially if the implant is deeply buried, migrated, or surrounded by dense, mature scar tissue, prolonging operative time and increasing risk.
  \item Allergic reaction to retained metal fragments: Though rare, this can occur if small pieces of metal are inadvertently left in situ, leading to chronic inflammation.
  \item Need for further surgery: To address complications such as refracture, infection, incomplete removal, or persistent symptoms.
\end{itemize}

\section*{Postoperative Recovery Timeline}
Immediately after metalwork removal, the patient is transferred to the Post-Anesthesia Care Unit (PACU) for close monitoring of vital signs, pain levels, and neurovascular status of the affected limb. Pain management is initiated, typically with oral or intravenous analgesics. Most metalwork removal procedures are performed on an outpatient basis, meaning the patient can go home the same day, provided they meet discharge criteria, have adequate pain control, and are ambulating safely. For more extensive removals, complex cases, or patients with significant comorbidities, an overnight hospital stay may be necessary for observation.

During the first 24-48 hours, patients are advised to rest, keep the surgical site clean and dry, and monitor for signs of infection such as excessive redness, swelling, warmth, or drainage. Swelling and bruising around the incision are common and typically resolve within one to two weeks. Activity restrictions are usually less stringent than after the initial fixation surgery, but patients are often advised to avoid strenuous activities, heavy lifting, or direct impact on the surgical site for several weeks to allow soft tissue healing and minimize the risk of refracture, particularly if large implants were removed from weight-bearing bones. Gradual return to full activity is guided by the surgeon, with full recovery and resolution of symptoms often occurring within 4-8 weeks, though complete remodeling of bone around screw holes can take several months.

\section*{Surgical Findings \& Pathology}
During metalwork removal, the surgeon meticulously documents several key findings. This includes the condition of the overlying soft tissues (e.g., presence of bursitis, inflammation, nerve impingement), the ease or difficulty of implant exposure due to scar tissue, and the condition of the bone surrounding the implant (e.g., complete union, stress shielding, osteolysis around screws, presence of screw hole defects). Any unexpected findings, such as signs of infection, implant breakage, or unusual tissue reactions, are also noted. The removed hardware itself is typically inspected for integrity and any signs of wear or corrosion.

Pathology reports are generally not generated for routine metalwork removal unless tissue samples were taken due to suspicion of infection, allergic reaction, or other unusual findings. If an infection is suspected, tissue cultures from the surgical site and/or a sample of periprosthetic tissue may be sent for microbiological analysis. If an allergic reaction is suspected, tissue around the implant may be sent for histological examination to look for specific inflammatory patterns. For most patients, the absence of the implant and the alleviation of symptoms are the key indicators of a positive result, and no formal pathology report on the hardware itself is typically issued.

\section*{Expected Postoperative Course}
A normal and successful postoperative course following metalwork removal is characterized by a progressive reduction in pain and discomfort, particularly the specific symptoms that prompted the surgery. Patients should experience uneventful wound healing, with the incision closing without signs of infection, excessive inflammation, or dehiscence. Initial post-operative pain, primarily from the surgical incision, is typically well-controlled with oral analgesics and diminishes significantly within the first week. Swelling and bruising around the surgical site are common but should gradually resolve. Patients are expected to regain their pre-injury or improved functional status, with restored range of motion and the ability to resume daily activities without implant-related discomfort. A gradual return to full activity, as guided by the surgeon, is anticipated, with the ultimate goal of enhancing the patient's quality of life by eliminating the issues caused by the retained hardware.

\section*{Postoperative Care \& Follow-up}
Following metalwork removal, patients typically have one or more follow-up appointments with their orthopedic surgeon.

\begin{itemize}
  \item \textbf{Initial Follow-up:} The first visit usually occurs 1-2 weeks post-operatively for a wound check, assessment of early healing, and removal of sutures or staples.
  \item \textbf{Subsequent Visits:} Subsequent visits may be scheduled at 4-6 weeks to evaluate functional recovery, assess pain levels, and provide guidance on activity progression and return to work or sport.
  \item \textbf{Imaging:} Routine imaging, such as X-rays, is generally not required unless there are specific concerns, such as persistent pain, suspicion of refracture, or incomplete hardware removal.
  \item \textbf{Rehabilitation:} Physical therapy or occupational therapy may be recommended, particularly if the patient experienced prolonged immobilization, has residual stiffness, muscle weakness, or requires specific guidance for regaining full function.
  \item \textbf{Activity Restrictions:} Patients are advised on a gradual return to full activity, with specific instructions tailored to the site of removal, the patient's bone quality, and the extent of the original injury.
  \item \textbf{Warning Signs:} Crucially, patients are educated on warning signs to promptly contact their doctor, including increasing pain, fever, chills, excessive redness, swelling, warmth, or purulent drainage from the incision site, which could indicate infection or other complications.
\end{itemize}

\section*{Epidemiology}
Metalwork removal is one of the most frequently performed elective orthopedic procedures globally. While precise incidence rates vary significantly based on the type of initial fixation, anatomical location, and patient population, estimates suggest that between 30\% and 50\% of all internal fixation devices are eventually removed. This procedure is common across all age groups, though the indications may differ; for instance, growth disturbance is a unique concern in pediatric patients, while pain and prominence are prevalent across all ages. There is no significant sex predilection for metalwork removal itself, as its epidemiology largely reflects that of the initial fractures or osteotomies. The most common indications for removal are pain or discomfort (reported in 40-60\% of removals), followed by palpable prominence or soft tissue irritation (20-30\%), and patient request. Infection, while a critical indication, accounts for a smaller percentage (typically 5-10\%) of elective removals. Mortality associated with elective metalwork removal is exceedingly low, reflecting its generally straightforward nature. Morbidity is primarily related to the procedure-specific risks outlined previously, with refracture rates varying from 0.5\% to 5\%, depending on the bone, implant type, and patient activity level.

\section*{Outcomes \& Prognosis}
The outcomes of metalwork removal are generally favorable, with high rates of patient satisfaction and symptom resolution. For patients undergoing removal due to pain or soft tissue irritation, success rates, defined as significant improvement or complete resolution of symptoms, typically range from 70\% to 90\%. Functional outcomes are also generally excellent, with most patients returning to their pre-injury activity levels or achieving improved function once the irritating hardware is removed. Recurrence of the original symptoms is uncommon if the hardware was indeed the primary cause of discomfort.

Prognostic factors for a successful outcome include a clear and direct correlation between the implant and the pre-operative symptoms, complete removal of all hardware components, and the absence of post-operative complications such as infection or refracture. While the procedure itself carries inherent risks, the long-term prognosis for patients who undergo metalwork removal for appropriate indications is overwhelmingly positive, leading to enhanced comfort, improved function, and a better quality of life. There are no specific landmark clinical trials like APPAC for appendicitis, but numerous cohort studies and systematic reviews consistently support the efficacy of symptomatic hardware removal.

\section*{References}
\begin{enumerate}
  \item MedlinePlus. Hardware removal. U.S. National Library of Medicine; 2025.
  \item Court-Brown CM, Heckman JD, McQueen MM, Ricci WM, Tornetta P. Rockwood and Green's Fractures in Adults. 9th ed. Wolters Kluwer; 2025.
  \item Canale ST, Beaty JH. Campbell's Operative Orthopaedics. 14th ed. Elsevier; 2025.
  \item American Academy of Orthopaedic Surgeons (AAOS). Management of Orthopaedic Implants. Clinical Practice Guideline; 2025.
\end{enumerate}

\clearpage